\chapterhead{3}{Write It So a Stranger Can Repeat It}{A test that cannot be repeated cannot be checked}\label{ch3}

Suppose you ran the paper clip test and found that the clip helped. A friend in another town reads your result. They want to know: is it true? The best answer you can give is not ``trust me''. It is a card with everything on it, so your friend can run the same test and see for themselves.

That is what scientists do. A result that nobody else can repeat is hard to build on. So the second half of a fair test is \textbf{writing down what you kept the same}.

\subhead{The test card}
Here is a card for the paper clip test. Notice that every line tells the reader something they would otherwise have to guess.

\begin{tcolorbox}[enhanced,colback=white,colframe=sky,boxrule=1.3pt,arc=2mm,left=8pt,right=8pt,top=5pt,bottom=5pt,before skip=7pt,after skip=7pt]
\sffamily\small
\textbf{Question.} Does a paper clip on the nose make a plane fly farther?\\[3pt]
\textbf{I change.} A paper clip: none, or one small clip on the nose.\\[3pt]
\textbf{I measure.} Distance from the launch line to where the plane first lands, nearest 0.1~m, using a tape measure.\\[3pt]
\textbf{I keep the same.}
\begin{itemize}[leftmargin=1.4em,itemsep=0pt,topsep=1pt]
\item Paper: A4 or letter sheets from the same pack.
\item Fold: the same dart fold, checked by laying the planes on top of each other.
\item Launch: from a marked line, standing still, arm at shoulder height, plane level.
\item Push: gentle, as if passing a cup of water. Same thrower for every throw.
\item Place: the same hall or yard, no wind, no obstacles.
\end{itemize}
\textbf{How many times.} Three throws each, in the mixed order plain, clip, plain, clip, plain, clip.\\[3pt]
\textbf{What I will report.} All six distances, the average of each group, and the spread.
\end{tcolorbox}

\subhead{Exact words beat vague words}
``Throw it hard'' means one thing to you and another to your friend. ``Same push'' is a start, but ``gentle, as if passing a cup of water'' is a picture the reader can copy. Whenever a line of the card says something vague, replace it with something someone else could repeat.

\begin{center}
\sffamily\small\renewcommand{\arraystretch}{1.5}
\begin{tabular}{@{}p{4.8cm}p{6.1cm}@{}}
\toprule
\textbf{Vague} & \textbf{Better}\\
\midrule
a normal piece of paper & A4 or letter paper from the same pack\\
throw it from a good spot & from behind a taped line, standing still\\
throw it the same & arm at shoulder height, gentle, plane level\\
measure how far it went & nearest 0.1~m from the line to where it first lands\\
do it a few times & three throws each, in a mixed order\\
\bottomrule
\end{tabular}
\end{center}

\begin{multix}[hmm]\multixbot[hmm]\tcblower
I have written a test card that says ``Throw the plane a good way.'' I think it is excellent. Anyone can follow it. Can they?
\end{multix}

\begin{tryit}{Write your own card, then trade}
Choose a question you can test with things at home. For example: ``Does a ramp that is higher make a toy car roll farther?'' Fill in the card below. Then give it to a friend or grown-up \emph{without talking}. Can they do your test using only the card?\par\medskip
\sffamily\small\renewcommand{\arraystretch}{2}
\begin{tabular}{@{}p{2.6cm}p{8.6cm}@{}}
\textbf{Question} & \wline[8.4cm]\\
\textbf{I change} & \wline[8.4cm]\\
\textbf{I measure} & \wline[8.4cm]\\
\textbf{I keep the same} & \wline[8.4cm]\\
 & \wline[8.4cm]\\
\textbf{How many times} & \wline[8.4cm]\\
\end{tabular}
\end{tryit}

\clearpage

\subhead{Did the stranger agree?}
Omar follows Mia's card exactly. He does not know what she found. Here are his practice numbers, a made-up run.

\begin{worked}{Omar repeats the test with Mia's card (practice numbers)}
\textbf{No clip:} 3.9, 4.3, 4.0~m. Total $3.9+4.3+4.0=12.2$. Average $12.2\div3\approx4.1$~m.\\
\textbf{Clip:} 5.0, 5.4, 5.5~m. Total $5.0+5.4+5.5=15.9$. Average $15.9\div3=5.3$~m.\\
Mia's averages were 4.1~m and 5.3~m. Omar's are also about 4.1~m and 5.3~m. The result was \textbf{repeated}.
\end{worked}

\begin{center}
\begin{tikzpicture}
  \node[bigcir,minimum size=2.5cm,fill=sky!8] (m) at (0,0) {};
  \node[cir,font=\sffamily\scriptsize] at (0,0.5) {4.1};\node[cir,font=\sffamily\scriptsize] at (0,-0.4) {5.3};
  \node[note] at (0,-1.55) {Mia: plain, clip};
  \node[bigcir,minimum size=2.5cm,fill=sky!8] (o) at (4.4,0) {};
  \node[cir,font=\sffamily\scriptsize] at (4.4,0.5) {4.1};\node[cir,font=\sffamily\scriptsize] at (4.4,-0.4) {5.3};
  \node[note] at (4.4,-1.55) {Omar: plain, clip};
  \draw[popA] (m) -- node[above,note] {same card} (o);
  \node[note,anchor=west,text width=4cm,align=left] at (6.1,0) {Two people, two sets of throws, one card. The averages match.};
\end{tikzpicture}
\end{center}

\trust{Another person followed the card and threw the planes themselves. Their averages, 4.1~m and 5.3~m, matched the first person's.}
{When a stranger gets about the same answer from the same written conditions, the result is not just one person's luck or habit. Every line of the card was clear enough to follow.}
{Whether the clip helps other planes, in other rooms, or with other paper. It also does not show \emph{why}. Matching averages show that the test can be repeated, not that the explanation is right.}

\begin{keepbox}
A good card lets someone else repeat your test. \textbf{We kept} \wline[3.4cm] \textbf{on the card, and we left off} \wline[3.4cm].
\end{keepbox}

\clearpage

\begin{puzzle}{The missing line}
Read this test card and find what is missing.
\begin{center}\sans\small
\begin{tabular}{|p{9cm}|}
\hline
\textbf{Question:} Does the length of a ramp change how far a toy car rolls?\\
\textbf{I change:} the length of the ramp, 30~cm or 60~cm.\\
\textbf{I measure:} how far the car rolls after the ramp, nearest cm.\\
\textbf{I keep the same:} the car, the floor, the person who lets go.\\
\hline
\end{tabular}
\end{center}
What one important thing is missing from ``I keep the same''? \wline[4.5cm]\\
Why would it matter? \wline[5.5cm]
\end{puzzle}

\begin{puzzle}{Put the steps in order}
These steps for the paper clip test are in a jumbled order. Write 1 to 6 beside them.
\begin{center}\sans\small\renewcommand{\arraystretch}{1.5}
\begin{tabular}{@{}lc@{}}
Throw the plain plane, then the clip plane, in a mixed order & \wline[1cm]\\
Fold two matching planes & \wline[1cm]\\
Find the average for each group & \wline[1cm]\\
Write down the distance after each throw & \wline[1cm]\\
Put a clip on one plane & \wline[1cm]\\
Tape a launch line on the floor & \wline[1cm]
\end{tabular}
\end{center}
\end{puzzle}

\begin{puzzle}{Make it exact}
Rewrite each vague line so that a stranger could follow it.
\begin{center}\sans\small\renewcommand{\arraystretch}{1.9}
\begin{tabular}{@{}p{4.3cm}p{6.2cm}@{}}
(a) ``Use a big cup.'' & \wline[6cm]\\
(b) ``Wait a while.'' & \wline[6cm]\\
(c) ``Stir it well.'' & \wline[6cm]
\end{tabular}
\end{center}
\end{puzzle}

\begin{puzzle}{Two people, two answers}
Two classes followed the same card for the clip test. Class A found that the clip plane flew 1.2~m farther. Class B found it flew only 0.1~m farther.\\
Write two things that could have been \emph{missing} from the card, so the classes did slightly different tests.\\
\rule{\linewidth}{0.5pt}\\[4pt]\rule{\linewidth}{0.5pt}
\end{puzzle}

\begin{puzzle}{Turn a story into a card}
Maya wrote: ``I rolled my car down a ramp, then I stacked more books under the ramp and rolled it again. It went farther.'' Fill in the card so a stranger could repeat her test.
\begin{center}\sans\small\renewcommand{\arraystretch}{1.9}
\begin{tabular}{@{}p{3cm}p{7.6cm}@{}}
Question & \wline[7.4cm]\\
I change & \wline[7.4cm]\\
I measure & \wline[7.4cm]\\
I keep the same & \wline[7.4cm]\\
How many times & \wline[7.4cm]
\end{tabular}
\end{center}
\end{puzzle}

\begin{puzzle}{Same card, different people}
Three people followed the same paper clip card. Their averages for the plain plane were 4.1, 4.0 and 6.9~m.\\
(a) Which person's result is different from the others? \wline[2.4cm]\\
(b) Write two questions you would ask that person about how they did the test.\\
\rule{\linewidth}{0.5pt}\\[4pt]\rule{\linewidth}{0.5pt}
\end{puzzle}

\begin{tryit}{Swap cards and compare}
Give your card from the earlier activity to a friend or grown-up. Ask them to run your test. Write your two averages and theirs below. Do the results agree?\par\medskip
\renewcommand{\arraystretch}{1.7}
\begin{tabular}{|l|c|c|}
\hline
 & \textbf{Version 1 average} & \textbf{Version 2 average}\\ \hline
\textbf{My results} & & \\ \hline
\textbf{Their results} & & \\ \hline
\end{tabular}\par\medskip
What did the card leave out? \wline[6.5cm]
\end{tryit}
\clearpage

\subhead{A second card: the ice test}
Here is the card for the ice test from Chapter~5. Read it, then check that every line could be followed by a stranger.

\begin{tcolorbox}[enhanced,colback=white,colframe=sky,boxrule=1.3pt,arc=2mm,left=8pt,right=8pt,top=5pt,bottom=5pt,before skip=7pt,after skip=7pt]
\sffamily\small
\textbf{Question.} Which wrapping keeps an ice cube from melting longest?\\[3pt]
\textbf{I change.} The wrapping: none, one layer of foil, one paper towel, or one wool sock.\\[3pt]
\textbf{I measure.} Minutes until the cube has fully melted, checked every 5 minutes with a timer.\\[3pt]
\textbf{I keep the same.}
\begin{itemize}[leftmargin=1.4em,itemsep=0pt,topsep=1pt]
\item Ice: four cubes from the same tray, taken out together.
\item Saucer: four identical saucers, side by side on one table.
\item Place: away from sunlight and drafts, in the same room.
\item Start: all four wrapped before the timer starts.
\end{itemize}
\textbf{How many times.} Two runs, with new ice each time.\\[3pt]
\textbf{What I will report.} Both times, the average for each wrapping, and the spread.
\end{tcolorbox}

\begin{puzzle}{Two things the card still leaves out}
This card says nothing about how \emph{thick} each wrapping is, or how many layers of the sock or towel to use. Write one more line for ``I keep the same'' that would fix it.\\
\rule{\linewidth}{0.5pt}\\[4pt]
Write one more thing a stranger might wonder about. \wline[5cm]
\end{puzzle}
\clearpage
